本文在 Shelley 叙事长诗 \emph{The Witch of Atlas} 上，用 100 题、带诗节级证据的 gold set 和 11 个同环境条件，检验了针对 LightRAG 的四个批评。四项判定并不一致。

\begin{itemize}
\item \textbf{归因（H1）得到支持。}在 28 个 chunk 的语料上，完整 LightRAG（0.845）相对 NaiveRAG（0.820）只高 0.025（$p=0.515$）；抽空 LLM 画像（0.840）或打乱一跳结构（0.830）也没有显著下降；去掉原始 chunk（0.445）则损失 0.400。正确率的主要来源是原文检索，而不是图结构或 LLM 改写。
\item \textbf{改写偏差（H2）只得到部分支持。}LLM 画像的自动支持率并不低于抽取式画像（0.952/0.871 对 0.883/0.877），答案层的引用精确率与 unsupported rate 也无显著差异；LLM 画像与原文的 ROUGE-L 更低（0.020/0.026 对 0.030/0.055），并存在 Hermaphroditus 这样的个案错误。“字面漂移”成立，“事实偏差”未被统计证实。
\item \textbf{时序（H3）只得到部分支持。}标准图在时序题上并不弱（0.88），“图普遍无法表示时序”不成立；图相对无图有时序定向收益（0.88 对 0.68，$p=0.005$），时序图在 20 题上达到 1.00（相对 C2 仅边缘显著，$p=0.077$），仅按时序重排 chunk（0.68）无效。时序结构有效，其上限依赖人工事件标注。
\item \textbf{评测与部署（H4）得到支持。}44 条人工盲审与评判模型的二次加权 $\kappa=0.797$，严格 $\kappa=0.535$；评判模型平均 1.57 对人工 1.30，系统性偏宽松。全生命周期成本为索引 32.5k token/195.8 秒、查询约 3.6k token/4.5 秒 P50，远高于“检索阶段 <100 token、一次调用”的口径；增量更新节省 50.9\%。
\end{itemize}

本文并不主张 LightRAG 无效。在小规模且原文可完整检索的语料上，图结构与 LLM 画像文本未表现出独立于原文的可检出增益；时序图存在方向一致但依赖标注的收益；评测与成本口径会显著影响结论。对实践者而言，先建立 chunk 基线、为原文保留充足预算、优先采用可追溯画像，并在时序任务中显式建模事件，比默认采用图结构更为稳妥。

本文的局限包括使用单一文学语料、7B 本地模型、单标注者、20 题时序子集与人工时间线等。未来工作应在百万 token 级的技术或法律语料上重复同一套 2$\times$2 归因设计，采用多标注者校准 judge，并将时序事件抽取自动化，以检验 C9 在真实部署场景中的性能上限。